% ==========================================
% ФАЙЛ: p10.tex
% БЛОК БИЛЕТОВ 46-50 (Степенные ряды и аналитичность)
% ==========================================

% --- Вопрос 46 ---
\examquestion{46}{Почленное интегрирование и дифференцирование рядов}
{Сформулируйте теоремы о почленном интегрировании и дифференцировании функциональных рядов. Докажите теорему о почленном интегрировании.}

\paragraph{Суть на пальцах}
Можно ли поменять местами знак суммы $\sum$ и интеграл $\int$ (или производную)? Да, но только если ряд сходится \textbf{равномерно}. Равномерная сходимость — это клей, который не дает бесконечной сумме развалиться при применении к ней операторов. Интеграл «добрый» (ему достаточно равномерности самого ряда), а производная «жесткая» (требует, чтобы равномерно сходился ряд из уже продифференцированных кусков).

\paragraph{1. Формулировки теорем}
Пусть функции $u_n(x)$ непрерывно дифференцируемы на отрезке $[a, b]$.
\begin{itemize}
    \item \textbf{Теорема об интегрировании:} Если ряд $\sum_{n=1}^{\infty} u_n(x)$ сходится \textbf{равномерно} на $[a, b]$ к сумме $S(x)$, то его можно интегрировать почленно:
    $$ \int_a^x \left( \sum_{n=1}^{\infty} u_n(t) \right) dt = \sum_{n=1}^{\infty} \left( \int_a^x u_n(t) dt \right) $$
    \item \textbf{Теорема о дифференцировании:} Если сам ряд $\sum u_n(x)$ сходится хотя бы в одной точке, а \textbf{ряд из производных} $\sum u'_n(x)$ сходится \textbf{равномерно} на $[a, b]$ к некоторой функции $\sigma(x)$, то исходный ряд сходится равномерно к $S(x)$, и его можно дифференцировать почленно: $S'(x) = \sigma(x)$.
\end{itemize}

\paragraph{2. Доказательство теоремы об интегрировании}
\begin{proof}
Представим сумму ряда в виде $S(x) = S_N(x) + R_N(x)$, где $S_N(x)$ — $N$-я частичная сумма, а $R_N(x)$ — остаток. Интегрируем обе части от $a$ до $x$:
$$ \int_a^x S(t) dt = \int_a^x S_N(t) dt + \int_a^x R_N(t) dt $$
Первый интеграл справа — это интеграл от \textbf{конечной} суммы. С ним проблем нет, интеграл суммы равен сумме интегралов: $\int S_N(t) dt = \sum_{n=1}^N \int u_n(t) dt$.

Теперь нужно доказать, что интеграл от остатка $\int_a^x R_N(t) dt$ стремится к нулю при $N \to \infty$. 
Так как ряд сходится \textbf{равномерно}, для любого $\varepsilon > 0$ найдется номер $N$, что для всех $t \in [a, b]$ выполняется $|R_N(t)| < \frac{\varepsilon}{b-a}$.
Оценим модуль интеграла от остатка:
$$ \left| \int_a^x R_N(t) dt \right| \le \int_a^x |R_N(t)| dt \le \int_a^b \frac{\varepsilon}{b-a} dt = \frac{\varepsilon}{b-a} \cdot (b-a) = \varepsilon $$
Так как $\int R_N \to 0$, переходя к пределу при $N \to \infty$, получаем:
$$ \int_a^x S(t) dt = \lim_{N \to \infty} \sum_{n=1}^N \int_a^x u_n(t) dt = \sum_{n=1}^{\infty} \int_a^x u_n(t) dt $$
Что и требовалось доказать.
\end{proof}
% --- Вопрос 46 ---
\begin{tcolorbox}[colback=gray!5, colframe=black, boxrule=0.5pt, arc=0pt]
\textbf{Резюмируем Вопрос 46 (Интегрирование и дифференцирование):} \\
\textbf{Тезис:} Равномерная сходимость позволяет менять местами операторы $\sum$ и $\int$ (или $d/dx$). \\
\textbf{Ход:} Для интегрирования $\int_a^x S(t)dt$ разбиваем сумму на конечную $S_N$ и остаток $R_N$. Интеграл от конечной суммы равен сумме интегралов. \\
\textbf{Трюк:} Оцениваем интеграл остатка: $|\int R_N dt| \le \int |R_N| dt \le \int \frac{\varepsilon}{b-a} dt = \varepsilon$ благодаря равномерной сходимости. \\
\textbf{Финал:} Интеграл остатка $\to 0$, значит предел суммы интегралов равен интегралу суммы. Для производной требования жестче: равномерно сходиться должен именно ряд производных.
\end{tcolorbox}



% --- Вопрос 47 ---
\examquestion{47}{Степенной ряд и первая теорема Абеля}
{Сформулируйте первую теорему Абеля. Проведите её доказательство методом мажорирования геометрической прогрессией. Обоснуйте существование радиуса сходимости.}

\paragraph{Суть на пальцах}
Степенной ряд $\sum a_n x^n$ — это бесконечный многочлен. Первая теорема Абеля говорит: если этот многочлен выжил (сошелся) в какой-то точке $x_0$, то он гарантированно, абсолютно и с запасом сходится во всех точках $x$, которые находятся \textbf{ближе} к нулю, чем $x_0$. Это создает строго симметричную «безопасную зону» сходимости вокруг нуля.

\paragraph{1. Первая теорема Абеля}
Если степенной ряд $\sum_{n=0}^{\infty} a_n x^n$ сходится в точке $x_0 \ne 0$, то он сходится \textbf{абсолютно} для любого $x$, удовлетворяющего условию $|x| < |x_0|$.

\paragraph{2. Доказательство}
\begin{proof}
Ряд $\sum a_n x_0^n$ сходится по условию. По необходимому признаку сходимости, его общий член стремится к нулю: $a_n x_0^n \to 0$. Раз последовательность имеет конечный предел, она ограничена. Значит, существует константа $M > 0$, такая что для всех $n$ выполнено:
$$ |a_n x_0^n| \le M $$

Возьмем любую точку $x$, такую что $|x| < |x_0|$. Преобразуем общий член ряда $\sum a_n x^n$ в этой точке, искусственно умножив и разделив на $x_0^n$:
$$ |a_n x^n| = \left| a_n x_0^n \cdot \left(\frac{x}{x_0}\right)^n \right| = |a_n x_0^n| \cdot \left| \frac{x}{x_0}\right|^n $$
Применим нашу оценку (заменим первый множитель на $M$):
$$ |a_n x^n| \le M \cdot q^n, \quad \text{где } q = \left| \frac{x}{x_0} \right| $$
Так как по условию $|x| < |x_0|$, то знаменатель дроби больше числителя, а значит $q < 1$. 
Справа мы получили члены бесконечно убывающей геометрической прогрессии $M \cdot q^n$. Ряд $\sum M q^n$ при $q<1$ сходится. Значит, по мажорантному признаку сравнения, исходный ряд из модулей $\sum |a_n x^n|$ сходится. Следовательно, ряд $\sum a_n x^n$ сходится \textbf{абсолютно}.
\end{proof}

\paragraph{3. Существование радиуса сходимости}
Из теоремы Абеля логически вытекает структура области сходимости. Множество всех $|x|$, при которых ряд сходится, ограничено нулем снизу. Пусть $R = \sup \{ |x| : \text{ряд сходится} \}$.
\begin{itemize}
    \item Если $|x| < R$, то по определению супремума найдется сходящаяся точка $x_0$ еще правее, и по теореме Абеля в $x$ ряд сойдется.
    \item Если $|x| > R$, ряд расходится (иначе супремум был бы больше).
\end{itemize}
Это число $R$ называется \textbf{радиусом сходимости}, а интервал $(-R, R)$ — интервалом сходимости.


\begin{tcolorbox}[colback=gray!5, colframe=black, boxrule=0.5pt, arc=0pt]
\textbf{Резюмируем Вопрос 47 (Первая теорема Абеля):} \\
\textbf{Тезис:} Если степенной ряд сходится в точке $x_0$, он сходится абсолютно везде «внутри» (при $|x| < |x_0|$). \\
\textbf{Ход:} Раз ряд сходится в $x_0$, его общий член $a_n x_0^n \to 0$ и ограничен константой $M$. \\
\textbf{Трюк:} В точке $x$ переписываем член как $|a_n x^n| = |a_n x_0^n| \cdot |x/x_0|^n \le M \cdot q^n$, где $q = |x/x_0| < 1$. \\
\textbf{Финал:} Ряд мажорируется сходящейся геометрической прогрессией. Это обосновывает симметричную структуру сходимости и существование точной границы — радиуса $R$.
\end{tcolorbox}




% --- Вопрос 48 ---
\examquestion{48}{Теорема Коши-Адамара}
{Сформулируйте и докажите теорему Коши-Адамара для вычисления радиуса сходимости степенного ряда. Проработайте доказательство в обе стороны.}

\paragraph{Суть на пальцах}
Теорема Коши-Адамара дает универсальную формулу для нахождения радиуса сходимости $R$ любого степенного ряда по его коэффициентам $a_n$. Она основана на радикальном признаке Коши (берем корень $n$-й степени). Из-за того, что коэффициенты могут вести себя хаотично (прыгать, обнуляться), мы обязаны использовать \textbf{верхний предел} ($\limsup$), чтобы поймать самую "тяжелую" подпоследовательность.

\paragraph{1. Формулировка теоремы Коши-Адамара}
Радиус сходимости $R$ степенного ряда $\sum_{n=0}^{\infty} a_n x^n$ вычисляется по формуле:
$$ \frac{1}{R} = \limsup_{n \to \infty} \sqrt[n]{|a_n|} $$
\textit{(Если предел равен $0$, то $R = \infty$; если предел равен $\infty$, то $R = 0$).}

\paragraph{2. Доказательство}
\begin{proof}
Применим радикальный признак Коши к ряду из модулей $\sum |a_n x^n|$. Вычислим верхний предел корня $n$-й степени из общего члена:
$$ L = \limsup_{n \to \infty} \sqrt[n]{|a_n x^n|} = |x| \cdot \limsup_{n \to \infty} \sqrt[n]{|a_n|} = \frac{|x|}{R} $$
Согласно признаку Коши, ряд сходится абсолютно, если $L < 1$, и расходится, если $L > 1$.

\textbf{Часть 1 (Сходимость внутри интервала $|x| < R$):}
Пусть $|x| < R$. Тогда $\frac{|x|}{R} < 1$, то есть $L < 1$. По признаку Коши ряд сходится абсолютно.

\textbf{Часть 2 (Расходимость вне интервала $|x| > R$):}
Пусть $|x| > R$. Тогда $L = \frac{|x|}{R} > 1$. Это означает, что верхний предел $\limsup \sqrt[n]{|a_n x^n|} > 1$.
По свойству верхнего предела это значит, что существует \textbf{подпоследовательность} номеров $n_k$, для которой члены $|a_{n_k} x^{n_k}| \ge 1$. 
Но если бесконечно много членов ряда по модулю больше единицы, то общий член ряда \textbf{не стремится к нулю} ($a_n x^n \not\to 0$). Нарушен необходимый признак сходимости. Ряд расходится.

\textbf{Вывод:} Число $R$, вычисленное по этой формуле, является точной границей между областью абсолютной сходимости и областью расходимости, то есть является радиусом сходимости степенного ряда.
\end{proof}
% --- Вопрос 48 ---
\begin{tcolorbox}[colback=gray!5, colframe=black, boxrule=0.5pt, arc=0pt]
\textbf{Резюмируем Вопрос 48 (Теорема Коши-Адамара):} \\
\textbf{Тезис:} Радиус сходимости $R$ вычисляется по формуле $1/R = \limsup_{n \to \infty} \sqrt[n]{|a_n|}$. \\
\textbf{Ход:} Применяем радикальный признак Коши к ряду из модулей $\sum |a_n x^n|$, получая предел $L = |x| \cdot \limsup \sqrt[n]{|a_n|} = |x|/R$. \\
\textbf{Трюк:} Если $|x| < R$, то $L < 1$ (абсолютная сходимость). Если $|x| > R$, то $L > 1$, что означает существование подпоследовательности, не стремящейся к нулю. \\
\textbf{Финал:} Из-за нарушения необходимого признака ($a_n x^n \not\to 0$) при $|x| > R$ ряд расходится. Формула точно разделяет зоны.
\end{tcolorbox}




% --- Вопрос 49 ---
\examquestion{49}{Свойства степенных рядов внутри интервала сходимости}
{Докажите, что степенной ряд сходится абсолютно и равномерно на любом замкнутом отрезке $[-r, r]$ внутри $(-R, R)$. Дайте обоснование свойствам непрерывности, интегрирования и дифференцирования.}

\paragraph{Суть на пальцах}
Интервал $(-R, R)$ — это открытая зона. Прямо у самих границ (в точках $-R$ и $R$) ряд может вести себя нестабильно. Но если мы отступим от краев хотя бы на миллиметр и возьмем любой замкнутый \textbf{внутренний} отрезок $[-r, r]$, ряд на нем станет идеальным: он «заморозится» (сойдется равномерно). А раз он равномерный, с ним можно делать всё: непрерывно рисовать, интегрировать и брать производные бесконечное число раз.

\paragraph{1. Доказательство равномерной сходимости на $[-r, r]$}
\begin{proof}
Пусть $R$ — радиус сходимости ряда $\sum a_n x^n$. Выберем любое число $r$ такое, что $0 < r < R$. 
Так как $r$ строго меньше радиуса сходимости, в точке $x = r$ ряд сходится \textbf{абсолютно}. Это значит, что числовой ряд $\sum |a_n r^n|$ сходится.

Рассмотрим функциональный ряд на отрезке $x \in [-r, r]$. Для любого $x$ из этого отрезка модуль его общего члена оценивается сверху:
$$ |a_n x^n| = |a_n| \cdot |x|^n \le |a_n| \cdot r^n = |a_n r^n| $$
Мы нашли для функционального ряда сходящуюся числовую мажоранту (крышу), не зависящую от $x$. По \textbf{мажорантному признаку Вейерштрасса}, исходный ряд сходится на отрезке $[-r, r]$ абсолютно и \textbf{равномерно}.
\end{proof}

\paragraph{2. Обоснование аналитических свойств}
Так как любой $x \in (-R, R)$ можно накрыть неким внутренним отрезком $[-r, r]$, на котором ряд сходится равномерно, применяются общие теоремы о функциональных рядах (Вопросы 45-46):
\begin{itemize}
    \item \textbf{Непрерывность:} Члены $a_n x^n$ непрерывны (это мономы). Ряд сходится равномерно. Значит, его сумма $S(x)$ \textbf{непрерывна} на всем $(-R, R)$.
    \item \textbf{Интегрирование:} Равномерная сходимость позволяет менять местами интеграл и сумму. Степенной ряд можно почленно интегрировать по любому отрезку внутри $(-R, R)$.
    \item \textbf{Дифференцирование:} Ряд из производных $\sum n a_n x^{n-1}$ имеет \textbf{тот же самый} радиус сходимости $R$ (так как $\sqrt[n]{n} \to 1$). Значит, продифференцированный ряд тоже сходится равномерно на $[-r, r]$. Следовательно, исходный ряд можно дифференцировать почленно. И так можно делать до бесконечности!
\end{itemize}

% --- Вопрос 49 ---
\begin{tcolorbox}[colback=gray!5, colframe=black, boxrule=0.5pt, arc=0pt]
\textbf{Резюмируем Вопрос 49 (Свойства внутри $(-R, R)$):} \\
\textbf{Тезис:} На любом внутреннем замкнутом отрезке $[-r, r]$ (где $r < R$) степенной ряд сходится абсолютно и равномерно. \\
\textbf{Ход:} В точке $r < R$ ряд сходится абсолютно. Для любого $x \in [-r, r]$ верно $|a_n x^n| \le |a_n r^n|$. \\
\textbf{Трюк:} Числовой ряд $\sum |a_n r^n|$ служит жесткой крышей. По признаку Вейерштрасса сходимость становится равномерной. \\
\textbf{Финал:} Равномерность разблокирует все аналитические свойства: непрерывность суммы, возможность почленного интегрирования и бесконечного дифференцирования.
\end{tcolorbox}



% --- Вопрос 50 ---
\examquestion{50}{Вещественные аналитические функции и классы гладкости}
{Сформулируйте определение аналитичности ($C^\omega$). Разберите пример Коши $e^{-1/x^2}$, докажите его бесконечную дифференцируемость и неаналитичность.}

\paragraph{Суть на пальцах}
Мы привыкли, что если функция гладкая (бесконечно дифференцируемая, класс $C^\infty$), то её можно разложить в ряд Тейлора. Оказывается, это \textbf{ложь}. Функция может быть идеально гладкой, но её ряд Тейлора будет равен нулю, хотя сама функция не ноль! Аналитические функции ($C^\omega$) — это элита: это только те функции, чей ряд Тейлора не просто существует, но и \textbf{сходится к самой функции}. 

\paragraph{1. Определения классов гладкости}
\begin{itemize}
    \item Функция $f(x) \in C^\infty(a, b)$, если она имеет непрерывные производные всех порядков на этом интервале.
    \item Функция $f(x)$ называется \textbf{аналитической} в точке $x_0$ (класс $C^\omega$), если в некоторой окрестности этой точки она может быть представлена сходящимся степенным рядом (рядом Тейлора): 
    $$ f(x) = \sum_{n=0}^{\infty} a_n (x - x_0)^n $$
\end{itemize}

\paragraph{2. Контрпример Коши («Ложная гладкость»)}
Рассмотрим функцию:
$$ f(x) = \begin{cases} e^{-1/x^2}, & x \ne 0 \\ 0, & x = 0 \end{cases} $$

\textbf{Шаг 1: Доказательство $f \in C^\infty$ и $f^{(n)}(0) = 0$.}
Найдем производные при $x \ne 0$. По правилу дифференцирования экспоненты:
$$ f'(x) = e^{-1/x^2} \cdot \frac{2}{x^3} $$
$$ f''(x) = e^{-1/x^2} \cdot \left( \frac{4}{x^6} - \frac{6}{x^4} \right) $$
Индукцией легко доказать, что любая $n$-я производная имеет вид:
$$ f^{(n)}(x) = e^{-1/x^2} \cdot P_k\left(\frac{1}{x}\right) $$
где $P_k$ — некий многочлен. 
Найдем производные в точке $x = 0$ по определению (через предел при $x \to 0$):
$$ f'(0) = \lim_{x \to 0} \frac{e^{-1/x^2} - 0}{x} = \lim_{t \to \infty} \frac{t}{e^{t^2}} = 0 \quad \text{(сделана замена } t = 1/x \text{)} $$
Экспонента в знаменателе растет катастрофически быстрее любого полинома $t^k$ в числителе (доказывается по Лопиталю). Поэтому предел для \textbf{любой} производной равен нулю:
$$ f^{(n)}(0) = \lim_{x \to 0} \frac{f^{(n-1)}(x)}{x} = 0 \quad \forall n \in \mathbb{N} $$
Значит, функция $f(x)$ бесконечно дифференцируема везде, включая $0$. $f(x) \in C^\infty$.

\textbf{Шаг 2: Доказательство неаналитичности.}
Попробуем построить для неё ряд Тейлора в точке $x_0 = 0$:
$$ \sum_{n=0}^{\infty} \frac{f^{(n)}(0)}{n!} x^n = 0 + 0 \cdot x + 0 \cdot x^2 + \dots \equiv 0 $$
Ряд Тейлора тождественно равен нулю! Но сама функция $f(x) = e^{-1/x^2}$ строго больше нуля в любой точке $x \ne 0$.
Следовательно, ряд Тейлора \textbf{не сходится к функции} ни в какой окрестности нуля. Функция \textbf{не является} аналитической в нуле.

% --- Вопрос 50 ---
\begin{tcolorbox}[colback=gray!5, colframe=black, boxrule=0.5pt, arc=0pt]
\textbf{Резюмируем Вопрос 50 (Пример Коши):} \\
\textbf{Тезис:} Бесконечной гладкости ($C^\infty$) недостаточно для аналитичности ($C^\omega$); ряд Тейлора может не сходиться к самой функции. \\
\textbf{Ход:} Исследуем $f(x) = e^{-1/x^2}$ (с $f(0)=0$). Все её производные содержат множитель $e^{-1/x^2}$ в числителе и полиномы в знаменателе. \\
\textbf{Трюк:} При $x \to 0$ экспонента затухает быстрее, чем растет любой полином $\implies f^{(n)}(0) = 0$ для всех $n$. Функция $\in C^\infty$. \\
\textbf{Финал:} Ряд Тейлора в нуле тождественно равен $0$, но $f(x) \ne 0$ при $x \ne 0$. Следовательно, функция не аналитическая.
\end{tcolorbox}